\chapter{Introduction}
\label{ch: introduction}

\vspace{0.66cm}

\setlength{\beforeepigraphskip}{0pt}  % tighten space above
\setlength{\afterepigraphskip}{0pt}   % tighten space below

\setlength{\epigraphwidth}{0.95\textwidth}
\epigraph{%
  The causes and consequences of disaster are not defined by an autonomous natural order, nor are they inevitable. Rather, they are bound up in human history, shaped by human action and inactions.%
}{Remes and Horowitz, \textit{Critial Disaster Studies}, 2021}

\vspace{0.66cm}


\noindent
Catastrophes triggered by natural hazards have accompanied human societies since their origins, albeit the experience of natural disaster continuously mutate throughout time and space. The 1755 Lisbon earthquake is marked by historians as a breaking point in the development of risk and climate science, when a single disaster becomes a data point -- a scientific, administrative, and legal object of study. In then industrializing nations, the growing importance of the probabilistic world as instrument to explain climate variability and extremes, allowed the implementation of interdisciplinary disaster risk reduction strategies, a now days well established expert discipline\footnote{In the midst of Enlightenment, the 1755 Lisbon earthquake reanimated philosophical debates around natural disasters and their celestial origins, dramatically challenging the idea of divine benevolence or any metaphysic and moral justifications to the disaster - as in Voltaire's \textit{Candide} (1759). Rather than placing emphasis on the futility of the catastrophe, in a less pessimistic spirit Jean-Jacques Rousseau emphasized instead the political and societal responsibility in determining the scale of the natural catastrophe, noting: \textit{``[...] convenez, par exemple, que, si la nature n'avait point rassemblé là vingt mille maisons de six à sept étages, et que si les habitants de cette grande ville eussent été dispersés plus également, et plus légèrement logés, le dégât eût été beaucoup moindre, et peut-être nul."} (Letter from J.J. Rousseau to M. de Voltaire, 18th August 1756).}. In France, the history of precautionary planning for natural disaster dates back at least at Eugène Belgrand with \textit{La Seine, études hydrologiques} (1872), where he describes the first real-time flood monitoring and warning system for the river Seine in France. These systems were put to test with the 1910 Great Flood of Paris, later coined as the ``flood of modernity” exactly because it spurred further technological and governance innovations, establishing France's modern precedent for state responsibility over disaster recovery and prevention. Today, in Paris, one can still see bridges and buildings engraved with century-old markers of the 1910's floods, serving as lasting reminders of the impact that extreme hydrological events can cause.
%(Lettre de J. J. Rousseau à Monsieur de Voltaire. Le 18. aout 1756)}
% a scientific, administrative, and legal object of study.\citep{Quenet2005} 
%\citeyear{Belgrand1872}

Anthropogenic climate change is further amplifying countries' exposure to natural catastrophe risks, by affecting the frequency, intensity, and timing of extreme weather events. The fat-tailed distribution of damages caused by extreme weather events is becoming ``fatter",
thereby challenging the validity of standard cost–benefit analyses underpinning large-scale infrastructure investment appraisals, as highlighted since \citeauthor{weitzman2009modeling}’s (\citeyear{weitzman2009modeling}) influential Dismal Theorem. As of 2025, floods and droughts are singled out as France’s two most consequential natural hazards  and are expected to remain so in coming decades (\citeauthor{CCR2023}, \citeyear{CCR2023}; \citeyear{CCR2025BilanCatNat}). France’s supreme audit institution \citep{CourDesComptes2022} points to a persistent insurance protection gap from natural catastrophes, estimating that inadequate prevention to flood risk alone is costing France’s fiscal position billions in avoidable losses each year. This gap feeds through the financial system and macroeconomy. A recent Banque de France climate stress-test \citep{ACPR2024ClimateStressTest} finds that delaying carbon pricing until 2030 would sharply raise banks’ credit risks through falling collateral values and higher defaults, and push a third of French insurers below solvency thresholds, triggering over €10 billion in recapitalisation needs for major mutuals. France Stratégie \citep{pisani2023economic} projects that even a ``moderate” 2°C warming scenario could trim up to 7\% from France's GDP and require more than €150 billion in adaptation spending by 2050. At the EU level, debates on how to address climate-related natural catastrophe risk are increasingly prominent, with the \cite{ECB2024insurancegap} recently proposing an integrated approach to manage rising tail risks linked to the growing frequency and severity of extreme weather events. \\

Against this backdrop, the present paper consists of an empirical investigation of extreme flooding events in France and aims to trace their effects on labor market dynamics and firm performance in France, with a focus on labor demand (employment, hours), wages (average and hourly), and proxies for output (total turnover), differentiating effects by sector, firm structure, and geographic location. Utilizing a matched employer–employee dataset covering all formal sector workers in France from 2000 to 2020, I exploit exogenous variation in the timing of flood events that trigger a ministerial declaration of natural disaster to identify the causal effect of these shocks on establishment and worker outcomes in affected municipalities. Flood events are identified based on the place-based criteria used by France’s public–private natural disaster compensation scheme, the \textit{Garantie CatNat}, which has recorded ministerial declarations of natural disasters at the municipality level since 1982. France maintains a dense historical record of flood episodes that led to official natural catastrophe decrees (\textit{arrêtés de catastrophe naturelle}) under the Cat-Nat scheme. To address identification challenges posed by staggered treatment and isolate the impact of flood-induced shocks from other confounding factors, I implement an event-study difference-in-differences (DiD) design drawing on recent advances in the DiD literature. In particular, the Local-Projection DiD estimator proposed in \cite{dube2023local} allows me to recover interpretable estimates in a flexible and efficient manner. To disentangle granular establishment-level effects from broader aggregate impacts of idiosyncratic flood shocks, I conduct the analysis at both the establishment and municipality level. My specifications target the average treatment effect on the treated (ATT). Given the localized nature of flood events and their quasi-random nature conditional on geographical location, the ATT provides the most relevant estimand for my research question, as my treatment is not applicable or plausible for the full population of establishments or municipalities in France. To address potential omitted variable bias -- such as differences in land-use regulation or local infrastructure investment -- and to reduce the risk of confounding between flood exposure and outcomes, I propose and develop a long-run flood risk index based on high-resolution floodplain maps for Europe and stratify the analysis by its quintiles, ensuring comparisons occur within similar baseline exposure bands. This accounts for the fact that firms and workers may sort into or avoid flood-prone municipalities based on unobserved characteristics correlated with observed outcomes. Moreover, the flood risk index I develop serves as a key covariate in estimating a propensity score for being flooded. By restricting the sample to regions of common support, I improve covariate balance and precision in the causal estimates. I also construct ATT weights that further enhance covariate balance and strengthen the credibility of the identification strategy. This approach is complemented by targeted sample restrictions and the inclusion of vectors of flood histories that aim to make flood shocks as good as random in the sample comparing units with similar past ``dosage" of flooding, an approached inspired by \cite{patel2023floods}. My framework also allows for estimating the effects by treatment intensity, as I measure flood intensity by its duration in days. I classify floods lasting less than one day as low-intensity events, and those lasting between 1 and 22 days or more as high-intensity. Because large scale natural disasters cause greater damage and are more likely to affect
establishments’ labour demand and wages, this approach enables me to disentangle the effects of short- versus long-duration floods. 

I find that, on average, floods lead to persistent reductions in employment (around 2\%) and modest declines in average wages (approximately 1\%), particularly evident at the end-of-year reporting date (31.12). Hourly wages are less affected, suggesting that worker productivity conditional on hours worked remains relatively stable. The most pronounced negative effects are concentrated in the commerce, construction, and local services sectors — especially among small, mono-establishment firms with limited geographic diversification. In contrast, industrial and agricultural sectors appear more resilient to flood shocks. Financial outcomes, such as turnover, exhibit weaker and more volatile responses, likely due to measurement limitations, slower adjustment dynamics, and mitigating mechanisms such as insurance payouts or public transfers. While firm-level effects are clear and persistent, commune-level results are more muted, reflecting aggregation effects and local recovery dynamics — a pattern consistent with findings in \cite{leiter2009creative} and \cite{usman2024going}. Reduced-form estimates from this study quantify the current impact of extreme flooding events. In turn, these estimates can be used to calibrate structural equilibrium models, thereby narrowing Knightian uncertainty about future damages from climate extremes—as in \cite{bilal2023anticipating}, where reduced-form evidence is fed into Integrated Assessment Models (IAMs). My findings equip firms, insurers, and governments with the evidence needed to shift from reactive to proactive adaptation strategies, tailored to sectoral, geographic, and firm-specific characteristics. Ultimately, future damages will depend on dynamic spatial interactions and constraints, a point underscored by recent literature \citep{balboni2025spatial} that employs dynamic spatial general equilibrium models to capture intertemporal decisions and spatial linkages across locations. \\


\noindent\textbf{Related Literature} This paper adds to the strand of environmental economics that measures the economic incidence of, and evaluates at different scales, the costs associated with weather shocks, with a focus on channels behind losses (\citeauthor{kahn2005death}, \citeyear{kahn2005death}; 
\citeauthor{dell2012temperature}, \citeyear{dell2012temperature}; 
\citeauthor{deryugina2018economic}, \citeyear{deryugina2018economic}; 
\citeauthor{hansel2020climate}, \citeyear{hansel2020climate}; 
\citeauthor{bilal2023anticipating}, \citeyear{bilal2023anticipating}). It contributes to at least four literatures: (i) physical risk and firm performance, and firms’ climate adaptation; (ii) urban and spatial economics used to project future damages; (iii) the micro- and macroeconomics of insurance against extreme weather events; and (iv) environmental inequality and optimal climate adaptation policy (e.g., \citeauthor{boyce2007inequality}, \citeyear{boyce2007inequality}; \citeauthor{chancel2023climate}, \citeyear{chancel2023climate}). These four strands are briefly reviewed below, with a particular focus on flood impacts.

Natural disasters as shocks to firms, and the broader literature on physical risk and firm performance, document multiple adjustment margins. Flood-related physical risk can alter firms’ capital structure \citep{ginglinger2023climate}, location and the distribution of activity \citep{castro2022climate}, aggregate productivity \citep{kocornik2020flooded}, and entry, employment, and output \citep{jia2022expecting}; it affects capital retirement, total factor productivity, and the dispersion of revenue-based marginal products of capital (\citeauthor{erda2024disasters}, \citeyear{erda2024disasters}; \citeauthor{prieto2024floods}, \citeyear{prieto2024floods}). Building on the seminal contribution by \cite{leiter2009creative}, physical risk from flooding has been shown to transmit to bank risk (\citeauthor{poschl2022banks}, \citeyear{poschl2022banks}; \citeauthor{gallagher2024natural}, \citeyear{gallagher2024natural}),  as well as through inflation \citep{ficarra2025weathering}. Investor sentiment and market valuation also respond \citep{hong2019climate}. Supply chains, credit conditions, product-market outcomes, and market competitiveness are further implicated (\citeauthor{barrot2016input}, \citeyear{barrot2016input}; \citeauthor{balboni2023firm}, \citeyear{balboni2023firm}; \citeauthor{zappala2023sectoral}, \citeyear{zappala2023sectoral}). Related adjustments by households and local governments include residential mobility \citep{lethimillock} and migration \citep{fan2024gone}, municipal spending \citep{morvan2022municipalities}, sovereign debt \citep{dey2022surging}, welfare dynamics (\citeauthor{bilal2023anticipating}, \citeyear{bilal2023anticipating}; \citeauthor{ciscar2011physical}, \citeyear{ciscar2011physical}), and infrastructure investment \citep{balboni2025harm}.

Urban and spatial economics provide tools to measure aggregate impacts and to project future damages by embedding reallocation, mobility, and adaptation into spatial equilibrium and dynamic frameworks. These frameworks complement reduced-form evidence and support counterfactual analysis of rising temperatures and associated flood risk; see \cite{carleton2024adaptation} for a review and \cite{balboni2025spatial} for recent advances. Debates about discounting and the social valuation of future extreme events (see the prominent Stern–Nordhaus debate, still ongoing and revisited in \citeauthor{stern2023climate}, \citeyear{stern2023climate}), the role of government, place-based and industrial policy, and state regulation all frame these modeling choices. Behavioral myopia can also lead societies to underweight past disasters, complicating inference and policy design \citep{gallagher2014learning}.

The economics of insurance markets for natural catastrophes emphasizes how access to credit and risk-transfer mechanisms shape disaster impacts. Many high-income countries regulate catastrophe insurance markets, and a range of flood-insurance schemes exist. Comparative institutional evidence spans Europe \citep{bock2024flood} and France’s \textit{CatNat} regime \citep{grislain2018natural}, addressed in the next section, while recent work discusses reforms to homeowners’ insurance under climate change \citep{boomhower2024insurance} and the role of ordeals in selection markets \citep{wagner2022adaptation}, where a distinction between demand-side and supply-side interventions can be made. Demand-side policies may involve giving homeowners access to reliable risk information and requiring the purchase of natural-disaster insurance, though uptake and enforcement are difficult. Supply-side reforms may include compulsory market participation and measures to foster risk-sharing, such as capital accumulation, reinsurance, and improved access to capital markets.

Environmental inequality and optimal climate adaptation policy stress that vulnerability to flooding—like climate risk more generally—reflects history, infrastructure, and politics \citep{rentschler2022flood}. Despite more defenses and insurance, real flood losses have risen; global vulnerability to hydro-hazards has fallen six-fold since the 1980s, yet damages have not declined commensurately \citep{kreibich2022challenge}. Inequality in impacts is stark: most flood-vulnerable populations live in developing countries \citep{rentschler2022flood}, while per-capita emissions remain higher in rich nations \citep{gandhi2022adapting}. An illuminating case study examining ex-ante adaptation channels in the Netherlands and Bangladesh—two low-lying, high-risk countries—shows how differences in early warning systems, infrastructure, and community capacity drive divergent outcomes in extreme flood events. Not accounting for these ex-ante channels could misleadingly suggest that the Netherlands is inherently less vulnerable, when in fact its superior risk-management regime is the key differentiator \citep{tol2018economic}. Justice-oriented policies, including EU carbon-tax–financed transfers to poorer countries, can mitigate disparities. Small upward shifts in the upper tail of flood-loss distributions can dominate expected-value calculations and undermine conventional cost–benefit tests; non-stationarity from warming, sea-level rise, and extremes compounds these challenges. \\


%(e.g., \citep{grover2024firm}; \citep{barrot2016input})
%(e.g., Balboni; \cite{jia2022expecting})
%In its focus on flooding and firms, it also relates closely to recent methodological innovations in measuring flood impacts. 

%Identification has advanced via sample restrictions to isolate exogenous floods (Erda), flood-history vectors (Patel), matching (New Zealand studies), and staggered timing (Italian work, Clo), alongside LP–DiD and stacked DiD designs (\citep{fatica2024floods}; \citep{roth2020local}; \citep{}; \citep{usman2024going}; \citep{ficarra2025weathering}); see also \cite{indaco2021hurricanes}. I return to the empirical strategy—where CatNat disaster declarations are assumed exogenous to unobserved outcome trajectories—in Section~\ref{ch:empirical_strategy}.


% \noindent\textbf{Related Literature} This paper adds to the strand of the environmental economics literature that aims to measure the economic incidence and evaluate, at different scales, the costs associated with shocks caused by the weather, focusing on the potential channels behind losses (\citep{kahn2005death}, \citep{dell2012temperature}; \citep{deryugina2018economic}, \citep{hansel2020climate}, \citep{bilal2023anticipating}). With its focus on the impact of flooding and on firms and workers adjustment margins this paper contributes to the following four strands of the literature: natural disasters as shocks to firms, physical risk and firm performance, and adaptation to climate change (e.g., \citep{grover2024firm}, \citep{barrot2016input}), Urban and spatial economics to project future damamges (Balboni,..), as well as the micro- and macroeconomics of insurance against extreme weather events (e.g.,  Chelle 2025, Carleton et al ), ad Environmental inequality and optimal climate adaptation policy ((e.g., \citep{boyce2007inequality}, \citep{chancel2023climate}), In its focus on flooding and firms it relates closely to recent methodologcial innovatiions (Pateel, Hussain, Erda) in measuring the impact of flooding. This section reviews these fours strands of the literarture, with a partocullar focus on flooding. The intitutional backgirn oF France in these regards is presented in the folllowing section. 

% Erda applies sample restrictions to isolate truly exogenous floods. Patel builds vectors of flood histories; New Zealand studies use matching,
% while Italian work (Clo) relies on staggered timing.
% Fatica et al.\ and several other papers employ LP-DiD or related stacked DiD designs; see the notes for a complete list. \cite{indaco2021hurricanes}.



% \textbf{Physical risk, firm performance, and firm's climate adatation} From a firm's perspective, the shock incurred by an extreme weather events and the consequences of climate-related physical risk may temporarily or permanently alter the return to critical production inputs such as labor, commercial structures or capital, imposing operational decisions to firms and effectively driving climate adaptation actions. Physical climate risk related to flooding has for example shown to alter the capital structure of firms \citep{ginglinger2023climate}, firms' location choices and distribution of economic activity \citep{castro2022climate} and aggregate productivity   \cite{kocornik2020flooded} , firm entry, emplolyment, and output \citep{jia2022expecting}), capital retirement and total factor productivity vs. value added (Phd Student discussed),  dispersison of revenue-based marginal product of capital orr other productivity measures \cite{zappala2023sectoral}, \citep{erda2024disasters}, \cite{prieto2024floods}, Leiter et al seminal contributioin..., transmition risk to bank risk in Denmak \cite{poschl2022banks} \cite{gallagher2024natural}, investors’ sentiment and market valuation (Hong et al. (2019), Credit conditions, the degree of competitivness in a market, supply chains as well as product market outcomes (\citep{barrot2016input}, \citep{balboni2023firm}), the degree of competition in markets (\citep{conteduca2024natural}). Meanwhile, individuals and local governments too suffer the consequences of such risk and adjusemnts mechanisms include residential mobility (\citep{lethimillock}), municipal spending (\citep{morvan2022municipalities}), sovereign debt (Bolton, 2022), changes in real wages, welfare dynamics (Bilal, SImialr to Bbilal at welfare level there is (JC Ciscar, A Iglesias, L Feyen, L Szabó, 2011), for Europe!), infrasteucture inivedtmetns (\citep{balboni2025harm}), consumption in the US, insurance premiums, housing. Back fo firms, recovering the impacts of natural disasters both at the microeconomic level and more aggregate level is important as aggregate data may help meausre impacts on aggregate welfare or economic actiivity and suggest that regions recover or even benefit from natural disasters, at a more granular, microeconomic level substantia negative impacts hidden by realocation effects across impacted and non-impacted firms in the same region result in positive aggregate dynamics, reallocation effects. \textbf{This highlights the importance of complementing macroeconomic models with granular data on exposure to natural catastrophes to fully capture the distributional and short- to medium-term consequences of climate-related hazards and inform the targeting of climate adaptation policies.}

% \textbf{Methodological advances in quantification and prediction of damages}: Debatets aroound climate change and how to future of the world from an onmic potn of vie dae abck at leas aththe faomus Stern vs Nordhaus conotrovery, still presetn today (Sigliltz). They focused on the corect discocunt factor of extreme evnts in the future siglitz on externalilty. coasebargain does not function and private optimum mmay differ, hence the externality that arises also because up front costs but long run returns (Draghi, inddustrial pocliy), hence the logic for govvernaemntal inttervention or public-private partnerhsips, and indeed mostly state regualted see OECD. 
% Secondy, disocunting too heavy of past events (Ghallagher), Behavioural myopia leads societies to underweight past disasters \cite{gallagher2014learning}.
% The random variation exploited to estimated causal relationship is not of political nature rather climate-related. Common sources of variation come from the temperature and quality of air, rainfall, see Even lava bombs! It is important to distinguish betweeen estimating current damages to firms and workers from flood disruptions, as well as to estimate the aggregate impacts and project future damages. To measure aggregate impacts or project current damages in the future, IAMs/structural models/dynamic spatial equilibrium models are often used, as they provide a way to test counterfactual scenarios about the development of rising temperatures and assciasted flood risk. For example, see spatial equilibrium model in \cite{balboni2025harm} \cite{jia2022expecting} etc..	These are reviewd in \cite{carleton2024adaptation}. An important topic is urban and spatial economics, which blends many elements together and is a rapidly expanding literature \citep{balboni2025spatial}. Secindly, advances in reduced form estimatation such as in Carleton shows improved an understanding of what iss meamsured with such reduced form estriatmes (i.e., ex-ante or also ex-post adaptation). SPecifically, stackedd regression that create groups with comparable ``dosage" of flooding have been recently used withintresting restulss (cite all.. \citep{ficarra2025weathering}, Dev Patel, Hussain, \citep{roth2020local}, \citep{fatica2024floods}, \citep{fan2024gone}, \citep{usman2024going}. It is challenging as ned to compare simialr ex-ante adpatation choices, and purge treatment of any endogeneity meaniing that is orthogonal to the error term. I will come back to the emprical strategy and identification asssumptions I meet in Section~\ref{ch: empirical_strategy}, where I will point out that my iodentification assumption is that  CatNat disaster declarations are exogenous with respect to any factors that affect the trajectories of our outcomes that are unobserved or
% omitted from the main equation. On the macro front; things like heterogenous agents models, ramsey model, Advanced Sustainability Economics (ETH Course sslildes), Science Po Syllabus Polyhenique. engodenoug growth models, cite these that les to new fields such as  IAMs. (connection IAM to classical macro).

% \textbf{The economics of insurance markets for natural catastrophes}: Access to credit and the existance of insurance mechanism are main determinants of the impact of disasters. The public healh response too to extreme events varies from country to country. Strain on social protection mechanisms, some countries do not have a budget. There exist many insurance mechanisms against floods (discuss Vox Dev Literature Review, and article on “Do flood insurance schemes in developing countries provide incentives to reduce physical risks?”). Government intervention in natural disaster insurance markets may be required. In many high income countries, the natural catastrophe insurance industry is indeed state-regulated\footnote{Demand side intervention/policy solutions include government agencies providing accurate risk information available to homeowners (role of sharing information rents); requiring homeowners to purchase natural disaster insurance (role of mandatory insurance, the challenge is the take-up/enforcement). On the other side, supply side intervention could involve mandating market participation, or reforming the the current regulatory environment to incentivize risk smoothing through either capital accumulation or reinsurance markets (reform could encourage participation of private insurers by facilitating access to capital markets that enable reinsurance and diversification) (NEEDS citation)}. AAn overview of different natural catstophe insurance regimes in Europe realted to flooding is laid out in \cite{bock2024flood}. The french natural catastorphe regime is specifically rviewd \cite{grislain2018natural} For a discussion of reforms to the homeowners insurance markets to account for climate change see Boomhower et al. (2024). Role of ordeals in selection markets. Discuss Wagner (2022), pag. 381. 

% \paragraph{Environmental inequality and optimal climate policy:}  

% cite hsiang disitrubioal clilamte change
% Vulnerability to flooding—like climate risk more generally—is a function of history, infrastructure, and politics \citep{rentschler2022flood, kreibich2022challenge}.

% Despite the proliferation of defences and insurance, flood losses have risen in real terms. Globally, vulnerability to hydro-hazards has fallen six-fold since the 1980s, yet damages have not declined commensurately \citep{kreibich2022challenge}. This lack of adequate coverage can be traced to several phenomena. First, inequality in climate impacts is well documented: most of the world’s flood-vulnerable populations live in developing countries \citep{rentschler2022flood}, even though per-capita emissions are far higher in rich nations \citep{gandhi2022adapting}. An illuminating and surprisingly powerful case study exaining ex-ante adaptation channels in the Netherlands and Bangladesh, two low-lying, high-risk countries, identifies how differences in early warning systems, infrastructure investments, and community capacity drive divergent outcomes in extreme flood events. Not accounting for these ex-ante channels could misleadingly lead one to conclude that the Netherlands is inherently less vulnerable, when in fact its superior risk management regime is the key differentiator \cite{tol2018economic}. Justice and new redistibution mechanisms can make that therre is less impact of climate change in poorer countries (Fabre), transfer to poorer contriess taken from EU carbon taxes. Explain in general (Mathhew gordon) that 



% a small upward shift in the upper tail in the distribution of flood losses in Europen countries can dominate established expected values and invalidate conventional cost–benefit tests. By rising tempetarures, sea rise, more extremme events become likelier, exacebating the vulnerabilites mmentioned earlier as these introduce non-stationary in the distributions of many natural systems and makes climate patterns no longer stable or predictable based on past trends, further complicate effective adaptation policy. However see in Indonesia coastal moral hazard in Indonesia \cite{hsiao2023sea} wherer governamental intervention can complicate things, or in Vietnam \cite{balboni2025harm} that consider the incentive of investing in infrsdtrutreu in coastal areas. Need incentives to make investmetnts that requies missed gains, oppounity costs of not doing that investments mamy be too big, need to rreduce that oppotunity cost. Whasinghton consesns, This goes back to Draghi. Exaplin that green indsutrial policy is neeeded and hohw my research could inform, also for equitable outcomes. 
% \href{https://www.ensae.fr/en/courses/6737}{Science Po Syllabus!} \\


\noindent The paper is organised as follows. Section \ref{ch: Background} explains the institutional background of flood and natural catastrophe declarations in France. Section \ref{ch: data} details the data used, and Section \ref{ch: empirical_strategy} the empirical strategy. Section \ref{ch: results} presents the results. Section \ref{ch: discussion} discusses mechanisms, policy implications, and links to the broader literature, while Section \ref{ch: conclusion} concludes.


